\par\Needspace{6\baselineskip}\medskip
\begin{center}\hypertarget{native-179}{}\hypertarget{native-179-article132}{}\pdfbookmark[2]{132.}{article-132}132.\end{center}
Statuendo brevitatis gratia
\[
\begin{gathered}
(0.1.2)=a,\qquad (0.\mathrm{I}.2)D'=-b,\\
(0.\mathrm{O}.2)D=c,\qquad (0.\mathrm{II}.2)D''=d,
\end{gathered}
\]
ita ut aequatio 10 art. 114 fiat
\[
a\delta'=b+c\cdot\frac{n}{n'}+d\cdot\frac{n''}{n'}\mathpunct{\raisebox{.45ex}{,}}
\]
coëfficientes $c$ et $d$ quidem erunt primi ordinis, facile vero ostendi potest,
differentiam $c-d$ ad secundum ordinem referendam esse. Hinc vero sequitur,
valorem quantitatis $\dfrac{cn+dn''}{n+n''}$ ex suppositione approximata $n:n''=\theta:\theta''$ prodeuntem
errore quarti tantum ordinis affectum esse, quin adeo quinti tantum, quoties
observatio media ab extremis aequalibus intervallis distat. Fit enim iste error
\[
=\frac{c\theta+d\theta''}{\theta+\theta''\mathpunct{\mkern-8mu\raisebox{-.65ex}{.}\mkern8mu}}
-\frac{cn+dn''}{n+n''}
=\frac{\theta\theta''(d-c)(\eta''-\eta)}{(\theta+\theta'')(\eta''\theta+\eta\theta'')}\mathpunct{\raisebox{.45ex}{,}}
\]
ubi denominator secundi ordinis est, numeratorisque factor alter $\theta\theta''(d-c)$
quarti, alter $\eta''-\eta$ secundi, vel in casu isto speciali tertii ordinis. Exhibita
itaque aequatione illa in hacce forma
\[
a\delta'=b+\frac{cn+dn''}{n+n''}\cdot\frac{n+n''}{n'}\mathpunct{\raisebox{.45ex}{,}}
\]
manifestum est, vitium methodi in art. praec. propositae non inde oriri, quod
quantitates $n,n''$ hisce $\theta,\theta''$ proportionales suppositae sunt, sed inde, quod
\textit{insuper} $n'$ ipsi $\theta'$ proportionalis statuta est. Hoc quippe modo loco factoris
$\dfrac{n+n''}{n'}$ valor minus exactus $\dfrac{\theta+\theta''}{\theta'}=1$ introducitur, a quo verus
\[
=1+\frac{\theta\theta''}{2\eta\eta''rr'r''\cos f\cos f'\cos f''}
\]
quantitate ordinis secundi discrepat (art. 128).
